%GregorioTeX common file (stuff needed by both gregoriotex and gregoriosyms)
%
% Copyright (C) 2007-2026 The Gregorio Project (see CONTRIBUTORS.md)
%
% This file is part of Gregorio.
%
% Gregorio is free software: you can redistribute it and/or modify
% it under the terms of the GNU General Public License as published by
% the Free Software Foundation, either version 3 of the License, or
% (at your option) any later version.
%
% Gregorio is distributed in the hope that it will be useful,
% but WITHOUT ANY WARRANTY; without even the implied warranty of
% MERCHANTABILITY or FITNESS FOR A PARTICULAR PURPOSE.  See the
% GNU General Public License for more details.
%
% You should have received a copy of the GNU General Public License
% along with Gregorio.  If not, see <http://www.gnu.org/licenses/>.

\gre@declarefileversion{gregoriotex-common.tex}{6.2.0}% GREGORIO_VERSION

\ifnum\luatexversion<100%
  \gre@error{Error: this document must be compiled with LuaTeX 1.0 or later}%
\fi%

% Run a command for each string in a comma-separated list of strings.
% #1: a comma-separated list of strings
% #2: a command
% Examples:
% \greparsecommas{}{\typeout} does nothing.
% \greparsecommas{foo}{\typeout} -> \typeout{foo}
% \greparsecommas{foo,bar}{\typeout} -> \typeout{foo}\typeout{bar}
% Note: Does not strip out any whitespace.
\newcount\gre@count@parsecommas@cur
\def\gre@parsecommas#1#2{%
  % https://tex.stackexchange.com/a/510410
  \def\gre@parsecommas@options{,#1,}%
  \gre@debugmsg{parsecommas}{\gre@parsecommas@options}%
  \IfStrEq{\gre@parsecommas@options}{,,}{%
    % do nothing
  }{%
    \StrCount{\gre@parsecommas@options}{,}[\gre@parsecommas@total]% number of items plus one
    \gre@count@parsecommas@cur=1\relax
    \loop\ifnum\gre@count@parsecommas@cur<\gre@parsecommas@total\relax
      \StrBetween[\the\gre@count@parsecommas@cur,\the\numexpr\gre@count@parsecommas@cur+1]{\gre@parsecommas@options}{,}{,}[\gre@parsecommas@option]%
      #2{\gre@parsecommas@option}%
      \advance\gre@count@parsecommas@cur by 1\relax
    \repeat
  }%
}

%%%%%%%%%
%% Debugging
%%%%%%%%%

\def\gre@debugmsg#1#2{%
  \IfStrEq{\gre@debug}{}{}%
    %false
    {\IfStrEq{#1}{all}%
      %true
      {\gre@error{Debug error: ‘all’ is not a permitted keyword}}%
      %false
      {\IfStrEq{\gre@debug}{,all,}%
        %true
        {\gre@typeout{GregorioTeX debug: (#1) #2}}%
        %false
        {\IfSubStr{\gre@debug}{,#1,}%
          %true
          {\gre@typeout{GregorioTeX debug: (#1) #2}}%
          %false
          {\relax}%
        }%
      }%
    }%
}%

\def\gre@fixdebug{%
  {% we're starting a group in order to localize \exploregroups
    \exploregroups%
    \StrRemoveBraces{\gre@debug}[\gre@debug]%
    \xdef\gre@debug{,\gre@debug,}%
  }%
}%

\IfStrEq{\gre@debug}{}{% then
}{% else
  \typeout{GregorioTeX is in debug mode}%
  \typeout{\gre@debug\space messages will be printed to the log.}%
  \gre@fixdebug\relax %
}%

%%%%%%%%%%%%%
% Macros to handle deprecation path.
%%%%%%%%%%%%%

% #1 - the deprecated macro
% #2 - the correct macro to use
\def\gre@deprecated#1#2{%
  \ifgre@allowdeprecated%
    \ifx&#2&%
      \gre@warning{#1\space is deprecated and marked for removal}%
    \else%
      \gre@warning{#1\space is deprecated.\MessageBreak Use #2\space instead}%
    \fi%
  \else%
    \ifx&#2&%
      \gre@error{#1\space is deprecated and marked for removal}%
    \else%
      \gre@error{#1\space is deprecated.\MessageBreak Use #2\space instead}%
    \fi%
  \fi%
  \relax%
}%

\def\gre@obsolete#1#2{%
  \ifx&#2&%
    \gre@error{#1\space is obsolete and no longer has any effect}%
  \else%
    \gre@error{#1\space is obsolete.\MessageBreak Use #2\space instead}%
  \fi%
  \relax%
}%

%%%%%%%%%%%%%
% Helper macro for visibility values
%%%%%%%%%%%%%

% Helper macro to process visibility values with phantom support and deprecation warnings
% #1 - the value to process (visible/invisible/phantom/hphantom)
% #2 - command to execute for visible (e.g., \gre@showlyricstrue)
% #3 - command to execute for invisible (e.g., \gre@showlyricsfalse)
% #4 - phantom wrapper count variable name (e.g., \gre@lyrics@phantomwrapper)
% #5 - macro name for error messages (e.g., \protect\gresetlyrics)
\def\gre@processvisibility#1#2#3#4#5{%
  \IfStrEqCase{#1}{%
    {visible}{#2#4=0\relax}%
    {invisible}{#3#4=0\relax}%
    {phantom}{#3#4=1\relax}%
    {hphantom}{#3#4=2\relax}%
  }[% all other cases
    \gre@error{Unrecognized option "#1" for #5\MessageBreak Possible options are: 'visible', 'invisible', 'phantom', and 'hphantom'}%
  ]%
}%

% Simplified version for commands that don't support phantom options
% #1 - option string (visible/invisible)
% #2 - command to execute for visible
% #3 - command to execute for invisible
% #4 - macro name for error messages
\def\gre@processvisibilitynophantom#1#2#3#4{%
  \IfStrEqCase{#1}{%
    {visible}{#2}%
    {invisible}{#3}%
  }[% all other cases
    \gre@error{Unrecognized option "#1" for #4\MessageBreak Possible options are: 'visible' and 'invisible'}%
  ]%
}%

% Simplified version for commands that accept visible/invisible/phantom with wrapper control
% #1 - option string (visible/invisible/phantom)
% #2 - command to execute for visible/phantom
% #3 - command to execute for invisible
% #4 - phantom wrapper count variable name
% #5 - macro name for error messages
\def\gre@processvisibilitywithphantom#1#2#3#4#5{%
  \IfStrEqCase{#1}{%
    {visible}{#2#4=0\relax}%
    {invisible}{#3#4=0\relax}%
    {phantom}{#2#4=1\relax}%
  }[% all other cases
    \gre@error{Unrecognized option "#1" for #5\MessageBreak Possible options are: 'visible', 'invisible', and 'phantom'}%
  ]%
}%

% Helper macro to apply phantom wrapper based on counter value
% #1 - phantom wrapper counter (0=no wrapper, 1=\phantom, 2=\hphantom)
% #2 - content to wrap
\def\gre@applyphantomwrapper#1#2{%
  \ifcase#1\relax
    % 0: no wrapper
    #2%
  \or
    % 1: \phantom
    \phantom{#2}%
  \or
    % 2: \hphantom
    \hphantom{#2}%
  \fi
}%

%%%%%%%%%%%%%
% Internal function tracing
%%%%%%%%%%%%%

% Each function in the internal name space (\gre@...) should have an invocation of
% \gre@trace as its first line with the argument being the name of the function (without
% the backslash) and the list of arguments.
% The last line of these functions should be \gre@trace@end.
% Macros which merely store a value are not considered to be functions and thus not
% included in the trace stack.
% We also exempt the internal aliases for functions defined elsewhere.
% The tracing messages are turned on via the debug interface, keyword: trace.

\def\gre@trace@prefix{GreTrace: }
\def\gre@trace#1{%
  \IfSubStr{\gre@debug}{,trace,}%
    {%print tracing messages and increase indent
      \gre@typeout{\gre@trace@prefix \unexpanded{#1}}%
      \edef\gre@trace@prefix{\gre@trace@prefix| }%
    }{}%do nothing if not tracing
}%
\def\gre@trace@end{%
  \IfSubStr{\gre@debug}{,trace,}%
    {%dedent trace messages and mark end
      \StrGobbleRight{\gre@trace@prefix}{2}[\gre@trace@prefix]%
      \gre@typeout{\gre@trace@prefix -----}%
    }{}%do nothing if not tracing
}

